% ============================================================================
%  Funciones hash y sus aplicaciones
%  Módulo 01 · Criptografía aplicada — ChainNotes
%
%  Versión LaTeX del apunte web temas/hash.html
%  Todas las figuras se generan con TikZ / pgfplots: no hay imágenes externas.
%
%  Compilar:  pdflatex funciones-hash.tex   (dos pasadas, por el índice)
%             o bien:  latexmk -pdf funciones-hash.tex
% ============================================================================

\documentclass[11pt,a4paper]{article}

% --- idioma y codificación -------------------------------------------------
\usepackage[T1]{fontenc}
\usepackage[utf8]{inputenc}
\usepackage[spanish,es-noshorthands,es-tabla]{babel}
\usepackage{lmodern}
\usepackage{microtype}

% --- página ----------------------------------------------------------------
\usepackage[a4paper,margin=2.6cm,headheight=14pt]{geometry}
\usepackage{fancyhdr}
\usepackage{enumitem}
\usepackage{booktabs}
\usepackage{tabularx}
\usepackage{array}
\usepackage{caption}

% --- matemáticas -----------------------------------------------------------
\usepackage{amsmath,amssymb,amsthm}
\usepackage{mathtools}

% --- gráficas --------------------------------------------------------------
\usepackage{tikz}
\usetikzlibrary{positioning,arrows.meta,calc,fit,backgrounds,shapes.geometric,
                decorations.pathreplacing}
\usepackage{pgfplots}
\pgfplotsset{compat=1.18}

% --- enlaces y cajas (hyperref antes que tcolorbox) ------------------------
\usepackage[colorlinks=true,linkcolor=violetdeep,urlcolor=violetdeep,
            citecolor=violetdeep,pdftitle={Funciones hash y sus aplicaciones},
            pdfauthor={ChainNotes}]{hyperref}
\usepackage[most]{tcolorbox}

% ---------------------------------------------------------------------------
%  Paleta (heredada de assets/css/styles.css)
% ---------------------------------------------------------------------------
\definecolor{violetc}{HTML}{7A68F5}
\definecolor{violetdeep}{HTML}{5B46E0}
\definecolor{mintc}{HTML}{1FA88A}
\definecolor{coralc}{HTML}{E04E5C}
\definecolor{amberc}{HTML}{C98A1F}
\definecolor{skyc}{HTML}{1E88C7}
\definecolor{inkc}{HTML}{282344}
\definecolor{bitone}{HTML}{3B3566}
\definecolor{bitzero}{HTML}{DCD9F0}

\hypersetup{linkcolor=violetdeep}

% ---------------------------------------------------------------------------
%  Estilo de párrafo y encabezados
% ---------------------------------------------------------------------------
\setlength{\parindent}{0pt}
\setlength{\parskip}{0.55em}
\captionsetup{font=small,labelfont={bf,color=violetdeep}}

% Este documento lleva muchas figuras grandes. Con los valores por omisión LaTeX
% las aparta a páginas propias y deja huecos; aflojamos los umbrales.
\renewcommand{\topfraction}{0.9}
\renewcommand{\bottomfraction}{0.8}
\renewcommand{\textfraction}{0.07}
\renewcommand{\floatpagefraction}{0.75}
\setcounter{topnumber}{3}
\setcounter{totalnumber}{4}

\pagestyle{fancy}
\fancyhf{}
\fancyhead[L]{\small\color{inkc}\textsc{Funciones hash y sus aplicaciones}}
\fancyhead[R]{\small\color{inkc}\thepage}
\renewcommand{\headrulewidth}{0.4pt}

% ---------------------------------------------------------------------------
%  Entornos de teoremas
% ---------------------------------------------------------------------------
\theoremstyle{definition}
\newtheorem{definicion}{Definición}
\theoremstyle{plain}
\newtheorem{teorema}{Teorema}
\newtheorem{proposicion}{Proposición}
\theoremstyle{remark}
\newtheorem*{observacion}{Observación}

% ---------------------------------------------------------------------------
%  Cajas de color
% ---------------------------------------------------------------------------
\newtcolorbox{nota}[1]{enhanced,breakable,colback=skyc!6,colframe=skyc!75!black,
  coltitle=white,fonttitle=\bfseries\small,title={#1},
  left=7pt,right=7pt,top=5pt,bottom=5pt,boxrule=0.7pt,arc=2.5pt}

\newtcolorbox{caso}[1]{enhanced,breakable,colback=amberc!8,colframe=amberc!85!black,
  coltitle=white,fonttitle=\bfseries\small,title={#1},
  left=7pt,right=7pt,top=5pt,bottom=5pt,boxrule=0.7pt,arc=2.5pt}

\newtcolorbox{aviso}[1]{enhanced,breakable,colback=coralc!7,colframe=coralc!85!black,
  coltitle=white,fonttitle=\bfseries\small,title={#1},
  left=7pt,right=7pt,top=5pt,bottom=5pt,boxrule=0.7pt,arc=2.5pt}

\newtcolorbox{laboratorio}[1]{enhanced,breakable,colback=violetc!6,colframe=violetc!80!black,
  coltitle=white,fonttitle=\bfseries\small,title={#1},
  left=7pt,right=7pt,top=5pt,bottom=5pt,boxrule=0.7pt,arc=2.5pt}

% ---------------------------------------------------------------------------
%  Atajos de notación
% ---------------------------------------------------------------------------
\newcommand{\bits}{\{0,1\}}
\newcommand{\dsha}{\mathrm{dSHA}}
\newcommand{\dH}{d_{\mathrm{H}}}
\newcommand{\concat}{\,\|\,}
\newcommand{\Adv}{\mathrm{Adv}}
\DeclareMathOperator{\HMAC}{HMAC}
\DeclareMathOperator{\KDF}{KDF}

% ---------------------------------------------------------------------------
%  Datos reales para la figura de avalancha.
%  Cada par a/b es el bit i-ésimo de SHA-256("blockchain") / SHA-256("Blockchain").
%    A = ef7797e13d3a75526946a3bcf00daec9fc9c9c4d51ddc7cc5df888f74dd434d1
%    B = 625da44e4eaf58d61cf048d168aa6f5e492dea166d8bb54ec06c30de07db57e1
%  Bits volteados: 130 de 256 (50.8 %).
% ---------------------------------------------------------------------------
\def\avalanchabits{%
  1/0, 1/1, 1/1, 0/0, 1/0, 1/0, 1/1, 1/0, 0/0, 1/1, 1/0, 1/1, 0/1, 1/1, 1/0, 1/1,%
  1/1, 0/0, 0/1, 1/0, 0/0, 1/1, 1/0, 1/0, 1/0, 1/1, 1/0, 0/0, 0/1, 0/1, 0/1, 1/0,%
  0/0, 0/1, 1/0, 1/0, 1/1, 1/1, 0/1, 1/0, 0/1, 0/0, 1/1, 1/0, 1/1, 0/1, 1/1, 0/1,%
  0/0, 1/1, 1/0, 1/1, 0/1, 1/0, 0/0, 1/0, 0/1, 1/1, 0/0, 1/1, 0/0, 0/1, 1/1, 0/0,%
  0/0, 1/0, 1/0, 0/1, 1/1, 0/1, 0/0, 1/0, 0/1, 1/1, 0/1, 0/1, 0/0, 1/0, 1/0, 0/0,%
  1/0, 0/1, 1/0, 0/0, 0/1, 0/0, 1/0, 1/0, 1/1, 0/1, 1/0, 1/1, 1/0, 1/0, 0/0, 0/1,%
  1/0, 1/1, 1/1, 1/0, 0/1, 0/0, 0/0, 0/0, 0/1, 0/0, 0/1, 0/0, 1/1, 1/0, 0/1, 1/0,%
  1/0, 0/1, 1/1, 0/0, 1/1, 1/1, 1/1, 0/1, 1/0, 1/1, 0/0, 0/1, 1/1, 0/1, 0/1, 1/0,%
  1/0, 1/1, 1/0, 1/0, 1/1, 1/0, 0/0, 0/1, 1/0, 0/0, 0/1, 1/0, 1/1, 1/1, 0/0, 0/1,%
  1/1, 0/1, 0/1, 1/0, 1/1, 1/0, 0/1, 0/0, 0/0, 1/0, 0/0, 0/1, 1/0, 1/1, 0/1, 1/0,%
  0/0, 1/1, 0/1, 1/0, 0/1, 0/1, 0/0, 1/1, 1/1, 1/0, 0/0, 1/0, 1/1, 1/0, 0/1, 1/1,%
  1/1, 1/0, 0/1, 0/1, 0/0, 1/1, 1/0, 1/1, 1/0, 1/1, 0/0, 0/0, 1/1, 1/1, 0/1, 0/0,%
  0/1, 1/1, 0/0, 1/0, 1/0, 1/0, 0/0, 1/0, 1/0, 1/1, 1/1, 1/0, 1/1, 0/1, 0/0, 0/0,%
  1/0, 0/0, 0/1, 0/1, 1/0, 0/0, 0/0, 0/0, 1/1, 1/1, 1/0, 1/1, 0/1, 1/1, 1/1, 1/0,%
  0/0, 1/0, 0/0, 0/0, 1/0, 1/1, 0/1, 1/1, 1/1, 1/1, 0/0, 1/1, 0/1, 1/0, 0/1, 0/1,%
  0/0, 0/1, 1/0, 1/1, 0/0, 1/1, 0/1, 0/1, 1/1, 1/1, 0/1, 1/0, 0/0, 0/0, 0/0, 1/1%
}

% ===========================================================================
\begin{document}
% ===========================================================================

% babel-spanish redefine \% inspeccionando \lastskip, lo que revienta dentro de
% los nodos de pgfplots. Restauramos la definición estándar de LaTeX: en este
% documento el espacio fino se escribe siempre a mano («50\,\%»).
\makeatletter
\DeclareRobustCommand{\%}{\char`\%\relax}
\makeatother

\begin{titlepage}
\centering
\vspace*{2.5cm}

{\Large\color{violetdeep}\textsc{ChainNotes}\par}
{\small\color{inkc}Módulo 01 · Criptografía aplicada\par}

\vspace{1.4cm}
{\Huge\bfseries Funciones hash\par}
\vspace{0.4cm}
{\LARGE y sus aplicaciones\par}

\vspace{1.2cm}
\begin{minipage}{0.78\textwidth}
\centering\itshape\color{inkc}
Una huella de 256 bits para cualquier cosa que exista: un nombre, una película,
una biblioteca entera. De ese único truco cuelgan la firma digital, el árbol de
Merkle, el compromiso, el minado y la inmutabilidad de la cadena. Aquí lo
desarmamos pieza por pieza.
\end{minipage}

\vspace{1.6cm}
\begin{tikzpicture}[x=6mm,y=6mm]
  \foreach \x in {0,...,7}{
    \pgfmathtruncatemacro{\t}{mod(\x*\x+3*\x+1,2)}
    \fill[violetc!\the\numexpr 25+10*\x\relax] (\x,0) rectangle ++(0.82,0.82);
  }
\end{tikzpicture}

\vfill
{\small Apuntes de Blockchain · Proyecto académico\par}
{\small\color{inkc}\today\par}
\end{titlepage}

\tableofcontents
\clearpage

% ===========================================================================
\section{Qué hace exactamente una función hash}
% ===========================================================================

\begin{nota}{Objetivo del módulo}
Reconocer las cuatro propiedades observables de una función hash
---determinismo, longitud fija, uniformidad y avalancha--- y distinguir un hash
criptográfico de uno que no lo es.
\end{nota}

Piensa en una licuadora industrial. Metes lo que sea ---una palabra, una
película, la Biblioteca Nacional--- y siempre sale exactamente un litro de puré.
Siempre un litro: ni más si metes más, ni menos si metes menos. Y el puré de una
manzana nunca se parece al puré de dos manzanas: no hay «casi igual».

\begin{figure}[htbp]
\centering
\begin{tikzpicture}[
  font=\small,
  ent/.style={draw=inkc!50,fill=white,rounded corners=2pt,minimum height=6.5mm,
              inner xsep=5pt,align=center},
  dig/.style={draw=violetc,fill=violetc!10,rounded corners=2pt,
              minimum height=6.5mm,minimum width=42mm,align=center},
  >={Stealth[length=2.2mm]}
]
% entradas
\node[ent,minimum width=14mm] (e1) at (0,1.5)  {\ttfamily "a"};
\node[ent,minimum width=30mm] (e2) at (0,0)    {\ttfamily contrato.pdf\,(4 MB)};
\node[ent,minimum width=44mm] (e3) at (0,-1.5) {\ttfamily biblioteca.tar\,(1.2 TB)};

% caja H
\node[draw=violetdeep,fill=violetc!15,rounded corners=3pt,minimum width=16mm,
      minimum height=32mm,align=center,font=\large] (H) at (5,0) {$H$};

% salidas
\node[dig] (s1) at (10.4,1.5)  {\ttfamily ca978112\ldots\,f8f0};
\node[dig] (s2) at (10.4,0)    {\ttfamily 9f86d081\ldots\,08a0};
\node[dig] (s3) at (10.4,-1.5) {\ttfamily 2c26b46b\ldots\,7ca7};

\foreach \i in {1,2,3}{
  \draw[->,inkc!60] (e\i.east) -- (H.west |- e\i);
  \draw[->,violetdeep] (H.east |- s\i) -- (s\i.west);
}

\node[align=center,font=\footnotesize,inkc] at (0,-2.8)
  {longitud arbitraria\\$\bits^{*}$};
\node[align=center,font=\footnotesize,violetdeep] at (10.4,-2.8)
  {longitud fija\\$\bits^{256}$};
\end{tikzpicture}
\caption{Una función hash aplasta un dominio infinito sobre un codominio finito.
Ésa es su utilidad y, a la vez, la razón por la que las colisiones son
inevitables.}
\label{fig:mapeo}
\end{figure}

\subsection{Tres cosas que la licuadora cumple}

\begin{description}[leftmargin=1.4em,style=unboxed,itemsep=0.3em]
  \item[Es determinista.] La misma manzana, el mismo puré, hoy y en diez años,
    en tu máquina y en la mía.
  \item[No se puede des-licuar.] Del puré no recuperas la manzana.
  \item[Es discontinua a propósito.] Quitarle un pixel a una foto no cambia «un
    poquito» el resultado: lo cambia por completo. Eso es el \emph{efecto
    avalancha}, y es la propiedad que vuelve útil a un hash para detectar
    alteraciones.
\end{description}

La figura~\ref{fig:rejilla} muestra esa última propiedad sin metáforas. Son los
256 bits de dos digests reales que difieren en una sola letra de la entrada. No
hay nada de mágico en ello: es el resultado buscado de sesenta y cuatro rondas de
sumas, rotaciones y XOR.

\begin{figure}[htbp]
\centering
\begin{tikzpicture}[x=3.6mm,y=3.6mm]
  % --- rejilla A ---
  \begin{scope}
    \foreach \a/\b [count=\k from 0] in \avalanchabits {
      \pgfmathtruncatemacro{\row}{int(\k/16)}
      \pgfmathtruncatemacro{\col}{mod(\k,16)}
      \ifnum\a=1 \fill[bitone] (\col,-\row) rectangle ++(0.86,0.86);
      \else      \fill[bitzero] (\col,-\row) rectangle ++(0.86,0.86);\fi
    }
    \node[font=\footnotesize\ttfamily,anchor=north west,inkc] at (0,-16.4)
      {SHA-256("blockchain")};
  \end{scope}
  % --- rejilla B ---
  \begin{scope}[xshift=79.2mm]
    \foreach \a/\b [count=\k from 0] in \avalanchabits {
      \pgfmathtruncatemacro{\row}{int(\k/16)}
      \pgfmathtruncatemacro{\col}{mod(\k,16)}
      \pgfmathtruncatemacro{\f}{abs(\a-\b)}
      \ifnum\b=1 \fill[bitone] (\col,-\row) rectangle ++(0.86,0.86);
      \else      \fill[bitzero] (\col,-\row) rectangle ++(0.86,0.86);\fi
      \ifnum\f=1 \draw[coralc,line width=0.45pt]
        (\col-0.07,-\row-0.07) rectangle ++(1.0,1.0);\fi
    }
    \node[font=\footnotesize\ttfamily,anchor=north west,inkc] at (0,-16.4)
      {SHA-256("Blockchain")};
  \end{scope}
  % leyenda
  \node[anchor=north,font=\footnotesize,inkc,align=center] at (18.9,-18.6)
    {\tikz\fill[bitone](0,0)rectangle(2.4mm,2.4mm);\,\ bit en 1 \qquad
     \tikz\fill[bitzero](0,0)rectangle(2.4mm,2.4mm);\,\ bit en 0 \qquad
     \tikz\draw[coralc,line width=0.45pt](0,0)rectangle(2.4mm,2.4mm);\,\ se volteó};
\end{tikzpicture}
\caption{Espejo de bits. Cambiar una mayúscula voltea \textbf{130 de los 256
bits}: el 50.8\,\%. Las celdas marcadas en rojo son las que cambiaron.}
\label{fig:rejilla}
\end{figure}

\subsection{La consecuencia incómoda: las colisiones existen}

Las entradas posibles son infinitas y las salidas son $2^{256}$. Por el
principio del palomar, hay infinitos pares de mensajes distintos con el mismo
hash. La criptografía no los elimina ---no puede---: sólo los vuelve
\emph{inencontrables}.

\begin{definicion}[Función hash criptográfica]
Una función hash es una aplicación
\[
  H : \bits^{*} \longrightarrow \bits^{n}
\]
de cadenas de longitud arbitraria a cadenas de longitud fija $n$, a la que se le
exige:
\begin{enumerate}[label=(\roman*),leftmargin=2.2em,itemsep=0.2em]
  \item \textbf{Eficiencia.} Calcular $H(x)$ cuesta poco: lineal en $|x|$ en la
    práctica.
  \item \textbf{Uniformidad.} Para $x$ al azar, $H(x)$ se distribuye
    aproximadamente uniforme sobre $\bits^{n}$. El modelo idealizado se llama
    \emph{oráculo aleatorio}.
  \item \textbf{Avalancha.} Si $x$ y $x'$ difieren en un bit, cada bit de
    $H(x)\oplus H(x')$ vale 1 con probabilidad $\approx \tfrac12$, de modo que
    \[
      \mathbb{E}\bigl[\dH(H(x),H(x'))\bigr] \;\approx\; \frac{n}{2}.
    \]
\end{enumerate}
\end{definicion}

\begin{laboratorio}{Laboratorio 1.1 — Medidor de avalancha}
El 50\,\% no basta, y CRC32 lo demuestra. El experimento toma 400 mensajes
aleatorios, les voltea siempre el mismo bit de entrada y mide la distancia de
Hamming entre los dos digests, expresada como fracción del digest.

\textbf{Qué observar.} CRC32 da un promedio de aspecto respetable, cerca del
50\,\%, y aun así produce \emph{un solo patrón de diferencia} con desviación
$\sigma = 0$: voltear ese bit cambia siempre exactamente los mismos bits de
salida, sin importar el mensaje. Eso no es mezcla, es una ecuación que el
atacante puede resolver. SHA-256 produce 400 patrones distintos en 400 pruebas y
una $\sigma$ cercana a $\sqrt{n}/2$. \emph{El promedio miente; los patrones no.}
\end{laboratorio}

\begin{figure}[htbp]
\centering
\begin{tikzpicture}
\begin{axis}[
  width=12.4cm, height=6.2cm,
  xlabel={fracción del digest que cambia, $\dH/n$},
  ylabel={pruebas (de 400)},
  xmin=0.37, xmax=0.63, ymin=0, ymax=440,
  xtick={0.40,0.45,0.50,0.55,0.60},
  xticklabel style={/pgf/number format/fixed,/pgf/number format/precision=2},
  grid=major, grid style={inkc!12},
  legend style={at={(0.02,0.97)},anchor=north west,font=\footnotesize,
                draw=inkc!30,fill=white},
  tick label style={font=\footnotesize}, label style={font=\small},
  axis lines=left
]
\addplot[ybar,bar width=5pt,fill=skyc!45,draw=skyc!80!black] coordinates {
  (0.3906,0.2) (0.4063,0.9) (0.4219,3.5) (0.4375,10.8) (0.4531,25.9)
  (0.4688,48.4) (0.4844,70.4) (0.5000,79.8) (0.5156,70.4) (0.5313,48.4)
  (0.5469,25.9) (0.5625,10.8) (0.5781,3.5) (0.5938,0.9) (0.6094,0.2)};
\addlegendentry{SHA-256 · 400 patrones distintos, $\sigma\approx\sqrt{n}/2$}

\addplot[ybar,bar width=5pt,fill=coralc!55,draw=coralc!80!black] coordinates {
  (0.4375,400)};
\addlegendentry{CRC32 · 1 solo patrón, $\sigma = 0$}

\draw[violetdeep,dashed,thick] (axis cs:0.5,0) -- (axis cs:0.5,440);
\node[violetdeep,font=\footnotesize,anchor=south west]
  at (axis cs:0.505,395) {$n/2$};
\end{axis}
\end{tikzpicture}
\caption{Distribución de la avalancha. Ambas funciones promedian cerca del
50\,\%; sólo una de las dos \emph{mezcla}. La campana es el sello de una moneda
lanzada por cada bit de salida; la barra única, el de una función lineal.}
\label{fig:avalancha}
\end{figure}

% ===========================================================================
\section{Las tres resistencias, y por qué no son la misma}
% ===========================================================================

Tres preguntas distintas, tres dificultades distintas. Lo único que las separa es
qué recibe el adversario y qué queda libre para que él lo elija.

\begin{figure}[htbp]
\centering
\begin{tikzpicture}[
  font=\small,
  panel/.style={draw=#1,fill=#1!5,rounded corners=3pt,minimum width=48mm,
                minimum height=52mm,anchor=north west},
  dado/.style={draw=inkc!55,fill=inkc!12,rounded corners=2pt,inner sep=3pt,
               font=\footnotesize},
  libre/.style={draw=inkc!45,dashed,fill=white,rounded corners=2pt,inner sep=3pt,
                font=\footnotesize},
  meta/.style={draw=#1,fill=#1!15,rounded corners=2pt,inner sep=3pt,
               font=\footnotesize},
  >={Stealth[length=1.8mm]}
]
% --- panel 1: preimagen
\node[panel=skyc] (p1) at (0,0) {};
\node[anchor=north west,font=\bfseries,skyc] at ($(p1.north west)+(3mm,-3mm)$)
  {Preimagen};
\node[anchor=north west,font=\scriptsize\itshape,inkc]
  at ($(p1.north west)+(3mm,-8mm)$) {one-way};
\node[dado] (a1) at ($(p1.north west)+(6mm,-17mm)$) {$y$ dado};
\node[meta=skyc] (b1) at ($(p1.north west)+(30mm,-17mm)$) {hallar $M$};
\draw[->,inkc] (a1) -- (b1);
\node[anchor=north west,align=left,font=\scriptsize,inkc,text width=40mm]
  at ($(p1.north west)+(3mm,-23mm)$)
  {Te doy una huella y te pido encontrar al dueño. El adversario no elige nada
   del enunciado.};
\node[anchor=south,font=\bfseries,skyc] at ($(p1.south)+(0,2mm)$) {$2^{n}$};

% --- panel 2: segunda preimagen
\node[panel=amberc] (p2) at (5.4,0) {};
\node[anchor=north west,font=\bfseries,amberc] at ($(p2.north west)+(3mm,-3mm)$)
  {2.\textsuperscript{a} preimagen};
\node[anchor=north west,font=\scriptsize\itshape,inkc]
  at ($(p2.north west)+(3mm,-8mm)$) {target collision};
\node[dado] (a2) at ($(p2.north west)+(6mm,-17mm)$) {$M$ dado};
\node[meta=amberc] (b2) at ($(p2.north west)+(31mm,-17mm)$) {$M'\neq M$};
\draw[->,inkc] (a2) -- (b2);
\node[anchor=north west,align=left,font=\scriptsize,inkc,text width=40mm]
  at ($(p2.north west)+(3mm,-23mm)$)
  {Te doy una persona concreta y te pido otra con su misma huella. El objetivo
   está fijado de antemano.};
\node[anchor=south,font=\bfseries,amberc] at ($(p2.south)+(0,2mm)$) {$2^{n}$};

% --- panel 3: colisión
\node[panel=coralc] (p3) at (10.8,0) {};
\node[anchor=north west,font=\bfseries,coralc] at ($(p3.north west)+(3mm,-3mm)$)
  {Colisión};
\node[anchor=north west,font=\scriptsize\itshape,inkc]
  at ($(p3.north west)+(3mm,-8mm)$) {collision-free · cumpleaños};
\node[libre] (a3) at ($(p3.north west)+(14mm,-14mm)$) {$M$ libre};
\node[libre] (a4) at ($(p3.north west)+(34mm,-14mm)$) {$M'$ libre};
\node[meta=coralc,font=\scriptsize] (b3)
  at ($(p3.north west)+(24mm,-23mm)$) {$H(M)=H(M')$};
\draw[->,inkc] (a3.south) -- (b3.north);
\draw[->,inkc] (a4.south) -- (b3.north);
\node[anchor=north west,align=left,font=\scriptsize,inkc,text width=42mm]
  at ($(p3.north west)+(3mm,-28mm)$)
  {Cualesquiera dos personas que compartan huella. Es la más fácil para el
   atacante: escoge ambos lados.};
\node[anchor=south,font=\bfseries,coralc] at ($(p3.south)+(0,2mm)$) {$2^{n/2}$};
\end{tikzpicture}
\caption{Las tres resistencias. Gris lleno $=$ lo que el adversario recibe;
contorno punteado $=$ lo que puede elegir. Cuanta más libertad, más barato el
ataque.}
\label{fig:resistencias}
\end{figure}

\begin{nota}{La consecuencia que hay que memorizar}
MD5 y SHA-1 están rotas \emph{en colisión}, no en preimagen. Siguen siendo
irreversibles; lo que ya no pueden hacer es garantizar que dos documentos
distintos tengan huellas distintas. Por eso una contraseña vieja en SHA-1 no se
recupera de golpe, pero una firma sobre SHA-1 no vale nada.
\end{nota}

\begin{proposicion}\label{prop:col-implica-2pre}
Si $H$ es resistente a colisiones, entonces es resistente a segunda preimagen.
\end{proposicion}

\begin{proof}[Demostración (por contrapuesta)]
Supongamos que existe un algoritmo eficiente $\mathcal{A}$ que, dado cualquier
$x_i$, devuelve $x_j \neq x_i$ con $H(x_i) = H(x_j)$. Elegimos $x_i$ al azar y
ejecutamos $\mathcal{A}(x_i)$. El par $(x_i, x_j)$ es, por definición, una
colisión hallada eficientemente. Luego $H$ no es resistente a colisiones.
\end{proof}

\begin{observacion}
Al romperse la resistencia a colisión todavía puede sobrevivir la de segunda
preimagen ---es el caso de SHA-1 hoy---, pero nunca al revés.
\end{observacion}

\begin{proposicion}\label{prop:reciproco}
El recíproco falla: existe $h$ resistente a colisiones que \emph{no} es
resistente a preimagen.
\end{proposicion}

\begin{proof}
Sea $g$ resistente a colisiones con salida de $n$ bits. Definimos
\[
  h(x) =
  \begin{cases}
    1 \concat x, & \text{si } |x| = n,\\[2pt]
    0 \concat g(x), & \text{en otro caso.}
  \end{cases}
\]
\emph{Es resistente a colisiones.} Si $h(x) = h(y)$, ambas salidas empiezan
igual. Si empiezan con 1 entonces $x = y$, contradicción; si empiezan con 0
entonces $g$ colisiona, contradicción.

\emph{No es resistente a preimagen.} Si $y$ empieza con 1, la preimagen se lee
directamente: es el resto de $y$.
\end{proof}

\begin{nota}{Moraleja de diseño}
Las tres propiedades deben pedirse por separado. La única implicación gratuita es
la de la Proposición~\ref{prop:col-implica-2pre}.
\end{nota}

\begin{laboratorio}{Laboratorio 2.1 — Caza de ataques sobre un hash truncado}
Con SHA-256 completo nadie encontrará una colisión, hoy ni en un millón de años
de cómputo. La seguridad no es magia: \textbf{es tamaño}. Truncando el digest a
$n$ bits el ataque cabe en un navegador.

\textbf{Qué observar.} La colisión aparece cerca de $2^{n/2}$ intentos; la
preimagen exige $2^{n}$. Con $n = 24$ son 4\,mil frente a 16 millones. Ésa es
toda la diferencia entre este módulo y el del cumpleaños.
\end{laboratorio}

% ===========================================================================
\section{Cómo está construida por dentro, y qué se rompe por eso}
% ===========================================================================

\begin{nota}{Objetivo del módulo}
Entender la construcción iterada de Damgård--Merkle y descubrir el ataque de
extensión de longitud, que es la razón de existir de HMAC.
\end{nota}

Un hash no procesa el mensaje completo de golpe: lo corta en bloques de 64 bytes
y los va digiriendo uno tras otro. Cada bloque se mezcla con el resultado del
anterior, como una masa que se amasa por tandas. El estado final, después del
último bloque, es el digest.

\begin{figure}[htbp]
\centering
\begin{tikzpicture}[
  font=\small,
  f/.style={draw=violetdeep,fill=violetc!12,rounded corners=2pt,
            minimum width=11mm,minimum height=9mm},
  blk/.style={draw=inkc!55,fill=white,rounded corners=2pt,
              minimum width=11mm,minimum height=6mm,font=\footnotesize},
  pad/.style={draw=amberc,fill=amberc!15,rounded corners=2pt,
              minimum width=11mm,minimum height=6mm,font=\footnotesize},
  >={Stealth[length=2mm]}
]
\node[font=\footnotesize,inkc] (iv) at (0,0) {$H_0 = \mathrm{IV}$};
\node[f] (f1) at (2.1,0) {$f$};
\node[f] (f2) at (4.7,0) {$f$};
\node[f] (f3) at (7.3,0) {$f$};
\node[f] (f4) at (9.9,0) {$f$};
\node[draw=mintc,fill=mintc!15,rounded corners=2pt,minimum width=11mm,
      minimum height=9mm] (g) at (12.3,0) {$g$};
\node[font=\footnotesize,violetdeep,anchor=west] (out) at (13.3,0) {$h(x)$};

\node[blk] (x1) at (2.1,1.6) {$x_1$};
\node[blk] (x2) at (4.7,1.6) {$x_2$};
\node[blk] (x3) at (7.3,1.6) {$\cdots$};
\node[pad] (x4) at (9.9,1.6) {$x_t\,\|\,|x|$};

\foreach \i in {1,2,3,4}{ \draw[->,inkc!70] (x\i) -- (f\i); }
\draw[->,inkc] (iv) -- (f1);
\draw[->,inkc] (f1) -- node[above,font=\scriptsize,inkc]{$H_1$} (f2);
\draw[->,inkc] (f2) -- node[above,font=\scriptsize,inkc]{$H_2$} (f3);
\draw[->,inkc] (f3) -- (f4);
\draw[->,inkc] (f4) -- node[above,font=\scriptsize,inkc]{$H_t$} (g);
\draw[->,violetdeep] (g) -- (out);

\node[font=\scriptsize,amberc,anchor=south,align=center] at (9.9,2.3)
  {relleno inyectivo\\(MD-strengthening)};
\end{tikzpicture}
\caption{Construcción iterada de Damgård--Merkle. El estado viaja de izquierda a
derecha; el digest \emph{es} el estado final.}
\label{fig:dm}
\end{figure}

\begin{teorema}[Damgård--Merkle]
Sea $f : \bits^{n}\times\bits^{b} \to \bits^{n}$ una función de compresión. Tras
aplicar un relleno inyectivo que codifica la longitud,
\[
  H_0 = \mathrm{IV}, \qquad
  H_i = f(H_{i-1}, x_i), \qquad
  h(x) = g(H_t).
\]
Si $f$ es resistente a colisiones, entonces $h$ también lo es.
\end{teorema}

\begin{proof}
Si los mensajes tienen distinto número de bloques, el relleno codifica longitudes
distintas en el último bloque, y ahí hay una colisión de $f$. Si tienen el mismo
número, se recorre la cadena hacia atrás: existe un primer índice donde las
entradas de $f$ difieren pero las salidas coinciden, es decir, una colisión de
$f$.
\end{proof}

\begin{aviso}{El teorema sólo transfiere colisión}
No dice nada sobre el comportamiento de $h$ como oráculo aleatorio, y ahí es
exactamente donde aparece la extensión de longitud.
\end{aviso}

\subsection{El ataque de extensión de longitud}

Esa arquitectura tiene una consecuencia que casi nadie ve venir: si conoces el
digest, \textbf{conoces el estado interno} con el que quedó la máquina. Y con el
estado en la mano puedes seguir amasando ---añadir texto al final del mensaje---
sin conocer el principio.

Un servidor autentica peticiones con $t = H(k \concat m)$, creyendo que sin la
clave nadie puede producir un $t$ válido. El atacante ve $m$ y $t$, nunca la
clave, y aun así puede calcular la etiqueta válida de
$m \concat \texttt{relleno} \concat z$ para el texto $z$ que se le antoje.

\begin{figure}[htbp]
\centering
\begin{tikzpicture}[
  font=\small,
  f/.style={draw=violetdeep,fill=violetc!12,rounded corners=2pt,
            minimum width=9mm,minimum height=8mm},
  fa/.style={draw=coralc,fill=coralc!12,rounded corners=2pt,
             minimum width=9mm,minimum height=8mm},
  blk/.style={draw=inkc!55,fill=white,rounded corners=2pt,inner sep=2.5pt,
              font=\footnotesize},
  >={Stealth[length=2mm]}
]
% fila legitima
\node[font=\footnotesize,inkc,anchor=east] (iv) at (0,1.4) {IV};
\node[f] (g1) at (1.3,1.4) {$f$};
\node[f] (g2) at (3.4,1.4) {$f$};
\node[f] (g3) at (5.5,1.4) {$f$};
\node[blk,violetdeep,anchor=west] (t) at (6.5,1.4) {$t = H(k\concat m)$};
\node[blk] at (1.3,2.6) {$k$};
\node[blk] at (3.4,2.6) {$m_1$};
\node[blk,draw=amberc,fill=amberc!15] at (5.5,2.6) {$m_2\|\text{pad}$};
\draw[->,inkc] (iv) -- (g1); \draw[->,inkc] (g1) -- (g2);
\draw[->,inkc] (g2) -- (g3); \draw[->,violetdeep] (g3) -- (t);
\node[anchor=east,font=\scriptsize,inkc,align=right] at (-0.2,1.4)
  {servidor\\(conoce $k$)};

% fila atacante
\node[blk,violetdeep] (t2) at (5.5,-0.6) {$t$};
\node[fa] (a1) at (7.8,-0.6) {$f$};
\node[fa] (a2) at (9.9,-0.6) {$f$};
\node[blk,coralc,anchor=west] (t3) at (10.9,-0.6) {$t'$ válida};
\node[blk,draw=coralc] at (7.8,0.6) {$z_1$};
\node[blk,draw=coralc] at (9.9,0.6) {$z_2\|\text{pad}$};
\draw[->,coralc,thick] (t2) -- node[above,font=\scriptsize,coralc]
  {reanuda el estado} (a1);
\draw[->,coralc] (a1) -- (a2); \draw[->,coralc] (a2) -- (t3);
\node[anchor=east,font=\scriptsize,coralc,align=right] at (4.6,-0.6)
  {atacante\\(nunca vio $k$)};

\draw[->,violetdeep,dashed] (t.south) .. controls +(0,-0.9) and +(0,0.9) ..
  (t2.north);
\end{tikzpicture}
\caption{Extensión de longitud. El digest publicado \emph{es} el estado interno:
el atacante retoma la cadena donde el servidor la dejó y autentica un mensaje que
nadie firmó.}
\label{fig:extension}
\end{figure}

\begin{caso}{Caso documentado}
En 2009 se falsificaron firmas de la API de Flickr exactamente así, sobre un MAC
artesanal de la forma $H(\text{secreto}\concat m)$. El error no estaba en SHA-1:
estaba en creer que un hash de Damgård--Merkle se comporta como un oráculo
aleatorio.
\end{caso}

\subsection{La construcción correcta: HMAC}

\[
  \HMAC_k(m) \;=\; H\bigl((k \oplus \mathrm{opad}) \concat
                    H((k \oplus \mathrm{ipad}) \concat m)\bigr)
  \tag{RFC 2104}
\]

El doble anidamiento hace que el estado interno final no sea el de un prefijo del
mensaje, y la extensión de longitud muere ahí. Es lo que usan TLS, JWT y la
derivación de monederos jerárquicos BIP32 ---con HMAC-SHA-512--- precisamente por
este motivo.

\begin{nota}{SHA-3 abandona esta arquitectura}
Keccak usa una construcción de \emph{esponja} con permutación
$\text{Keccak-}f[1600]$, absorción a razón $r$ y capacidad $c = 2n$. No sufre
extensión de longitud, y ésa fue una de las razones de la elección del NIST en
2012. Bitcoin resuelve el mismo problema por otra vía: aplica doble SHA-256 a los
encabezados de bloque.
\end{nota}

% ===========================================================================
\section{El ataque del cumpleaños: por qué 256 bits valen 128}
% ===========================================================================

En un salón con 23 personas la probabilidad de que dos compartan cumpleaños
supera el 50\,\%. La intuición falla porque uno se pregunta por sí mismo, cuando
la pregunta real es por \emph{todos los pares posibles}: con 23 personas hay 253
pares.

Traducido a hashes: no cuesta $2^{256}$ romper SHA-256 por colisión, cuesta
$2^{128}$. La raíz cuadrada es el precio de que el atacante pueda elegir ambos
mensajes.

\subsection{Derivación de la cota}

Sea $N = 2^{n}$ el número de salidas y $k$ el de mensajes hasheados al azar. La
probabilidad de no colisión es
\[
  \Pr[\text{sin colisión}] \;=\; \prod_{i=0}^{k-1}\frac{N-i}{N}.
\]
Usando $1 - t \le e^{-t}$ con $t = i/N$,
\[
  \Pr[\text{colisión}] \;\approx\; 1 - e^{-k(k-1)/2N}.
\]
Imponiendo probabilidad $\tfrac12$ y aproximando $k(k-1)\approx k^{2}$,
\begin{equation}\label{eq:cumple}
  k \;=\; \sqrt{2\ln 2}\,\sqrt{N} \;\approx\; 1.1774\cdot 2^{n/2}.
\end{equation}

\begin{figure}[htbp]
\centering
\begin{tikzpicture}
\begin{axis}[
  name=ax1, width=7.4cm, height=5.6cm,
  xlabel={personas $k$}, ylabel={Pr[colisión]},
  xmin=0,xmax=100, ymin=0,ymax=1.03,
  grid=major, grid style={inkc!12}, axis lines=left,
  tick label style={font=\footnotesize}, label style={font=\small},
  title={\small $N = 365$ días}, title style={font=\small,color=inkc}
]
\addplot[skyc,very thick,domain=0:100,samples=120]
  {1-exp(-x*(x-1)/(2*365))};
\addplot[only marks,mark=*,mark size=1.3pt,coralc] coordinates {
  (10,0.1169)(20,0.4114)(30,0.7063)(40,0.8912)(50,0.9704)(60,0.9941)(70,0.9992)};
\draw[violetdeep,dashed] (axis cs:23,0) -- (axis cs:23,0.507);
\draw[violetdeep,dashed] (axis cs:0,0.507) -- (axis cs:23,0.507);
\node[violetdeep,font=\scriptsize,anchor=west] at (axis cs:25,0.44)
  {$k=23 \Rightarrow 50.7\,\%$};
\end{axis}

\begin{axis}[
  at={(ax1.right of south east)}, anchor=left of south west,
  width=7.4cm, height=5.6cm,
  xlabel={$k/\sqrt{N}$ (muestras normalizadas)}, ylabel={Pr[colisión]},
  xmin=0,xmax=3.2, ymin=0,ymax=1.03,
  grid=major, grid style={inkc!12}, axis lines=left,
  tick label style={font=\footnotesize}, label style={font=\small},
  title={\small la misma curva, para cualquier $N$},
  title style={font=\small,color=inkc},
  legend style={at={(0.97,0.06)},anchor=south east,font=\scriptsize,
                draw=inkc!30,fill=white}
]
\addplot[violetc,very thick,domain=0:3.2,samples=120] {1-exp(-x*x/2)};
\addlegendentry{$1-e^{-u^{2}/2}$}
\addplot[only marks,mark=*,mark size=1.3pt,skyc] coordinates {
  (0.523,0.1169)(1.047,0.4114)(1.570,0.7063)(2.094,0.8912)(2.617,0.9704)};
\addlegendentry{$N=365$}
\addplot[only marks,mark=triangle*,mark size=1.6pt,coralc] coordinates {
  (0.391,0.0728)(0.781,0.2621)(1.172,0.4961)(1.563,0.7048)(1.953,0.8517)
  (2.344,0.9361)(2.734,0.9764)};
\addlegendentry{$N=2^{16}$}
\draw[amberc,dashed,thick] (axis cs:1.1774,0) -- (axis cs:1.1774,1.0);
\node[amberc,font=\scriptsize,anchor=south west,rotate=90]
  at (axis cs:1.24,0.04) {$k = 1.1774\sqrt{N}$};
\end{axis}
\end{tikzpicture}
\caption{La paradoja del cumpleaños es una sola curva. Al medir las muestras en
unidades de $\sqrt{N}$, el salón de clases y un espacio de $2^{16}$ digests caen
sobre la misma gráfica: por eso la ecuación~\eqref{eq:cumple} no depende del
contexto.}
\label{fig:cumpleanos}
\end{figure}

\begin{nota}{Consecuencia operativa}
Una función hash de $n$ bits tiene \emph{seguridad ideal} si hallar una colisión
cuesta $\Theta(2^{n/2})$ y hallar una preimagen cuesta $\Theta(2^{n})$. Todo
ataque publicado que baje de esa cota se considera una \textbf{rotura}, aunque
siga siendo impracticable. Es la única cifra que permite decidir si una noticia
de criptoanálisis importa o no.
\end{nota}

Ésta es la razón por la que ningún diseño moderno se conforma con 128 bits de
salida: 128 bits de digest son sólo 64 bits de seguridad frente a colisión, y
$2^{64}$ es alcanzable. Es exactamente la historia de MD5, y la
figura~\ref{fig:costos} lo pone en una sola imagen.

\begin{figure}[htbp]
\centering
\begin{tikzpicture}
\begin{axis}[
  width=12.4cm, height=6.2cm,
  xlabel={longitud de salida $n$ (bits)},
  ylabel={$\log_2$ del trabajo esperado},
  xmin=100,xmax=280, ymin=0,ymax=270,
  xtick={128,160,192,224,256}, ytick={0,32,64,96,128,192,256},
  grid=major, grid style={inkc!12}, axis lines=left,
  tick label style={font=\footnotesize}, label style={font=\small},
  legend style={at={(0.03,0.97)},anchor=north west,font=\footnotesize,
                draw=inkc!30,fill=white}
]
\fill[coralc!10] (axis cs:100,0) rectangle (axis cs:280,64);
\node[coralc,font=\scriptsize,anchor=west] at (axis cs:104,34)
  {zona alcanzable por un adversario bien financiado ($\le 2^{64}$)};
\addplot[skyc,very thick,domain=100:280,samples=2] {x};
\addlegendentry{preimagen: $2^{n}$}
\addplot[coralc,very thick,domain=100:280,samples=2] {x/2};
\addlegendentry{colisión: $2^{n/2}$}
\addplot[only marks,mark=*,mark size=2pt,inkc] coordinates {(128,64)(160,80)(256,128)};
\node[inkc,font=\scriptsize,anchor=north west] at (axis cs:128,62) {MD5};
\node[inkc,font=\scriptsize,anchor=north west] at (axis cs:160,78) {SHA-1};
\node[inkc,font=\scriptsize,anchor=south east] at (axis cs:254,130) {SHA-256};
\draw[inkc!50,dashed] (axis cs:100,64) -- (axis cs:280,64);
\end{axis}
\end{tikzpicture}
\caption{Colisión contra preimagen. La línea de colisión de MD5 nace justo sobre
la frontera de lo alcanzable: su ruptura no fue un accidente, era cuestión de
tiempo.}
\label{fig:costos}
\end{figure}

\begin{laboratorio}{Laboratorio 4.1 — La paradoja, en personas y en bits}
El mismo cálculo sirve para un salón de clases y para un espacio de digests.

\textbf{Qué observar.} Con 365 días y 23 personas la probabilidad exacta es
50.7\,\%, y la aproximación exponencial casi no se distingue de ella. Al subir el
espacio a $2^{24}$ hacen falta unas 4\,800 muestras, no 8 millones: la raíz
cuadrada, otra vez.
\end{laboratorio}

% ===========================================================================
\section{Cuáles usar hoy, cuáles ya cayeron, y cuándo}
% ===========================================================================

Ninguna función hash muere de golpe. Primero alguien publica un ataque que cuesta
un poco menos que la fuerza bruta ---nadie se alarma---, luego el costo baja lo
suficiente para un laboratorio, y un día aparecen dos archivos distintos con el
mismo hash, descargables por cualquiera. En ese momento la función ya llevaba
años muerta en los papers.

\begin{figure}[htbp]
\centering
\begin{tikzpicture}
\begin{axis}[
  xbar, width=11.4cm, height=6.4cm,
  bar width=7pt,
  xlabel={$\log_2$ del costo de la mejor colisión conocida},
  xmin=0, xmax=152, xtick={0,32,64,96,128},
  symbolic y coords={SHA3-256,SHA-256,SHA-1,SHA-0,MD5},
  ytick=data,
  grid=major, grid style={inkc!12},
  tick label style={font=\footnotesize}, label style={font=\small},
  legend style={at={(0.5,-0.20)},anchor=north,legend columns=2,
                font=\footnotesize,draw=inkc!30,fill=white,
                /tikz/every even column/.append style={column sep=12pt}},
  enlarge y limits=0.16,
  nodes near coords, nodes near coords align={horizontal},
  every node near coord/.append style={font=\scriptsize,color=inkc,
    /pgf/number format/set decimal separator={.}}
]
\addplot[fill=skyc!35,draw=skyc!80!black] coordinates {
  (64,MD5) (80,SHA-0) (80,SHA-1) (128,SHA-256) (128,SHA3-256)};
\addlegendentry{seguridad ideal $2^{n/2}$}
\addplot[fill=coralc!55,draw=coralc!80!black] coordinates {
  (24.1,MD5) (39,SHA-0) (63.1,SHA-1) (128,SHA-256) (128,SHA3-256)};
\addlegendentry{mejor colisión publicada}
\draw[coralc,dashed,thick] (axis cs:64,SHA3-256) -- (axis cs:64,MD5);
\end{axis}
\end{tikzpicture}
\caption{Línea del deterioro. La distancia entre las dos barras es la medida
objetiva del daño. SHA-256 y SHA3-256 no tienen rotura conocida, así que sus dos
barras coinciden. Estado del arte a agosto de 2026.}
\label{fig:deterioro}
\end{figure}

\begin{table}[htbp]
\centering\small
\caption{Qué usar, qué no, y por qué. La cifra que importa no es la longitud de
la salida, sino la \emph{longitud de seguridad}.}
\label{tab:algoritmos}
\begin{tabularx}{\textwidth}{@{}l l >{\raggedright\arraybackslash}X l@{}}
\toprule
\textbf{Algoritmo} & \textbf{Seguridad} & \textbf{Mejor ataque conocido} &
\textbf{Estado}\\
\midrule
SHA-256 \newline {\scriptsize FIPS 180-4} &
  128 colisión \newline 256 preimagen &
  Cumpleaños genérico. El criptoanálisis dedicado no rebasa 31 de las 64 rondas. &
  \textcolor{mintc}{\textbf{vigente}}\\
\addlinespace
SHA3-256 \newline {\scriptsize Keccak · FIPS 202} &
  128 colisión \newline 256 preimagen &
  No supera 6 de las 24 rondas de Keccak-$f$. La esponja es inmune a extensión de
  longitud. &
  \textcolor{mintc}{\textbf{vigente}}\\
\addlinespace
SHA-1 \newline {\scriptsize retirada por el NIST} &
  $\approx 63$ efectivos &
  Colisión de prefijo elegido: $2^{63.4}$, ejecutada en 2020 por
  $\approx 45\,000$ USD de GPU rentada. &
  \textcolor{coralc}{\textbf{rota}}\\
\addlinespace
MD5 \newline {\scriptsize RFC 1321} &
  $< 32$ efectivos &
  Colisión diferencial de Wang: segundos en una laptop. Con prefijo elegido,
  $\approx 2^{39}$. &
  \textcolor{coralc}{\textbf{rota}}\\
\addlinespace
CRC32 \newline {\scriptsize suma de comprobación} &
  0 &
  Invertible en tiempo lineal. Segunda preimagen y colisiones triviales. &
  \textcolor{inkc}{no cripto.}\\
\addlinespace
Argon2id \newline {\scriptsize PHC 2015 · RFC 9106} &
  parametrizable &
  Función de derivación de claves: memoria y tiempo configurables. No es para
  integridad. &
  \textcolor{violetdeep}{contraseñas}\\
\bottomrule
\end{tabularx}
\end{table}

\begin{aviso}{Dos precisiones que suelen citarse mal}
La cifra $2^{39}$ corresponde a \textbf{SHA-0}, no a SHA-1. Y SHA-3 fue
seleccionado en octubre de 2012 y publicado como FIPS 202 en agosto de 2015: no
son la misma fecha.
\end{aviso}

% ===========================================================================
\section{Contraseñas: por qué un hash rápido es un defecto}
% ===========================================================================

\begin{nota}{Objetivo del módulo}
Entender qué protege la sal, qué protege el factor de trabajo, y por qué SHA-256
es la elección equivocada para almacenar contraseñas.
\end{nota}

Guardar $\mathrm{SHA\text{-}256}(\text{contraseña})$ parece prudente: si roban la
base, no hay contraseñas en claro. Pero el atacante no necesita invertir el hash.
Le basta con hashear las diez mil contraseñas más comunes y buscar coincidencias
---y como el hash es determinista, \textbf{una sola tabla sirve para todos los
usuarios de todas las bases de datos del mundo}. Eso es una tabla arcoíris.

La \emph{sal} rompe esa reutilización: un valor aleatorio distinto por usuario
obliga a rehacer la tabla para cada uno. El \emph{factor de trabajo} ataca lo
otro: hace que cada intento cueste 100\,ms en lugar de $0.000002$\,s.

\subsection{Los tres escenarios}

\begin{enumerate}[leftmargin=2.2em,itemsep=0.35em]
  \item \textbf{Ingenuo:} $v = H(p)$. El adversario invierte $|D|$ evaluaciones
    una sola vez y ataca a $u$ usuarios en $O(|D| + u)$.
  \item \textbf{Con sal:} $v = H(s \concat p)$, con $s$ aleatoria pública de al
    menos 128 bits. El costo pasa a $O(u\cdot|D|)$: la precomputación deja de
    amortizarse.
  \item \textbf{Con hash adaptativo:} $v = \KDF(p, s, \tau)$. El defensor paga
    $c$ por autenticación; el atacante paga $|D|\cdot c$ \emph{por usuario}:
    \[
      \mathrm{costo}_{\text{atacante}} \;=\; u \cdot |D| \cdot c(\tau, m).
    \]
\end{enumerate}

\begin{figure}[htbp]
\centering
\begin{tikzpicture}
\begin{axis}[
  width=12.4cm, height=6.6cm,
  xmode=log, ymode=log,
  xlabel={usuarios en la base de datos, $u$},
  ylabel={evaluaciones equivalentes de SHA-256},
  xmin=1, xmax=1e7, ymin=1e6, ymax=1e22,
  grid=major, grid style={inkc!12}, axis lines=left,
  tick label style={font=\footnotesize}, label style={font=\small},
  legend style={at={(0.03,0.97)},anchor=north west,font=\footnotesize,
                draw=inkc!30,fill=white},
  legend cell align=left
]
\addplot[coralc,very thick,domain=1:1e7,samples=60] {1e7 + x};
\addlegendentry{ingenuo $H(p)$: $|D| + u$}
\addplot[amberc,very thick,domain=1:1e7,samples=60] {1e7*x};
\addlegendentry{con sal: $u\cdot|D|$}
\addplot[mintc,very thick,domain=1:1e7,samples=60] {1e7*x*1e5};
\addlegendentry{sal $+$ KDF ($c = 10^{5}$): $u\cdot|D|\cdot c$}
\draw[violetdeep,dashed,thick] (axis cs:1,1.8e19) -- (axis cs:1e7,1.8e19);
\node[violetdeep,font=\scriptsize,anchor=south east] at (axis cs:9e6,2.2e19)
  {$2^{64}$: frontera de lo practicable};
\end{axis}
\end{tikzpicture}
\caption{Costo total del ataque por diccionario ($|D| = 10^{7}$ candidatos). Sin
sal, el esfuerzo del adversario es prácticamente constante por más usuarios que
haya; la sal lo vuelve lineal, y el factor de trabajo lo empuja cinco órdenes de
magnitud hacia arriba.}
\label{fig:contrasenas}
\end{figure}

\begin{caso}{Caso documentado}
LinkedIn, 2012: 6.5 millones de hashes SHA-1 sin sal. La mayoría se recuperó en
días. De ahí salieron bcrypt, scrypt y Argon2.
\end{caso}

\begin{nota}{Criterio de diseño}
Ajusta $\tau$ para que una verificación legítima cueste entre 50 y 250\,ms en tu
hardware de producción. Ése es el punto en que el usuario no lo nota y el
atacante sí.
\end{nota}

\begin{laboratorio}{Laboratorio 6.1 — Tabla arcoíris contra sal}
Ocho usuarios ficticios eligen contraseñas de un diccionario común. Se ejecuta el
ataque con y sin sal, mirando la columna de \emph{hashes calculados}, no la de
cuentas rotas.

\textbf{Qué observar.} La sal no impide adivinar una contraseña débil concreta:
impide \emph{atacar a todos a la vez con el mismo trabajo}. Con ocho usuarios el
atacante paga ocho veces más; con ocho millones, ocho millones de veces más.
\end{laboratorio}

% ===========================================================================
\section{Árboles de Merkle: probar pertenencia sin descargar todo}
% ===========================================================================

Un bloque de Bitcoin puede contener miles de transacciones. Tu teléfono no va a
descargar 600\,GB para comprobar que te pagaron.

El árbol de Merkle resuelve exactamente eso: hashea las transacciones por
parejas, luego los resultados por parejas, y así hasta quedarse con un único
valor ---la \emph{raíz}--- que resume todo el conjunto.

\begin{definicion}[Construcción y prueba de inclusión]
Dadas $m$ hojas, el nivel base es
$L_0 = (\dsha(T_1),\dots,\dsha(T_m))$ con
$\dsha(x) = \mathrm{SHA256}(\mathrm{SHA256}(x))$, y recursivamente
\[
  L_{j+1}[i] \;=\; \dsha\bigl(L_j[2i] \concat L_j[2i+1]\bigr),
\]
duplicando el último elemento cuando el nivel tiene un número impar de nodos.

Una \emph{prueba de inclusión} para $T_i$ es la lista de hermanos
$\pi = (\sigma_0,\dots,\sigma_{d-1})$ con $d = \lceil \log_2 m \rceil$, y el
verificador recomputa
\[
  r_0 = \dsha(T_i), \qquad
  r_{j+1} = \dsha\bigl(\mathrm{orden}(r_j, \sigma_j)\bigr),
\]
aceptando si $r_d = \text{raíz}$.
\end{definicion}

\begin{figure}[htbp]
\centering
\begin{tikzpicture}[
  font=\footnotesize,
  nodo/.style={draw=inkc!45,fill=white,rounded corners=2pt,
               minimum width=13mm,minimum height=5.5mm,font=\scriptsize\ttfamily},
  camino/.style={draw=violetdeep,fill=violetc!35,line width=0.9pt},
  herm/.style={draw=amberc,fill=amberc!40,line width=0.9pt},
  raiz/.style={draw=mintc,fill=mintc!35,line width=1pt},
  hoja/.style={draw=inkc!35,fill=inkc!6,rounded corners=2pt,
               minimum width=13mm,minimum height=5mm,font=\scriptsize},
  >={Stealth[length=1.6mm]}
]
% nivel 0 (hojas)
\foreach \i/\h [count=\c from 0] in
  {1/f66c00,2/8f6903,3/ac7e1b,4/2ff4a3,5/e419d9,6/6135e7,7/113e7d,8/24a231}{
  \node[nodo] (L0\i) at ({\c*1.75},0) {\h};
  \node[hoja] (T\i) at ({\c*1.75},-0.9) {$T_{\i}$};
  \draw[->,inkc!50] (T\i) -- (L0\i);
}
\node[camino,nodo] at (L03) {ac7e1b};
\node[herm,nodo]   at (L04) {2ff4a3};

% nivel 1
\foreach \i/\h [count=\c from 0] in {1/77b663,2/9e7a54,3/cdf6fd,4/6bb813}{
  \node[nodo] (L1\i) at ({\c*3.5+0.875},1.5) {\h};
}
\node[herm,nodo]   at (L11) {77b663};
\node[camino,nodo] at (L12) {9e7a54};

% nivel 2
\foreach \i/\h [count=\c from 0] in {1/17a876,2/e58276}{
  \node[nodo] (L2\i) at ({\c*7+2.625},3) {\h};
}
\node[camino,nodo] at (L21) {17a876};
\node[herm,nodo]   at (L22) {e58276};

% raiz
\node[raiz,nodo,minimum width=16mm] (R) at (6.125,4.5) {5db548};
\node[mintc,font=\scriptsize\bfseries,anchor=west] at (7.2,4.5) {raíz de Merkle};

% aristas
\foreach \a/\b in {1/1,2/1,3/2,4/2,5/3,6/3,7/4,8/4}{
  \draw[inkc!45] (L0\a) -- (L1\b);
}
\foreach \a/\b in {1/1,2/1,3/2,4/2}{ \draw[inkc!45] (L1\a) -- (L2\b); }
\draw[inkc!45] (L21) -- (R); \draw[inkc!45] (L22) -- (R);

% resaltar camino
\draw[violetdeep,line width=1pt] (L03) -- (L12);
\draw[violetdeep,line width=1pt] (L12) -- (L21);
\draw[violetdeep,line width=1pt] (L21) -- (R);

% leyenda
\node[anchor=west,align=left,font=\scriptsize,inkc] at (-0.9,-2.2)
  {\tikz\node[camino,minimum width=4mm,minimum height=2.6mm]{}; camino que
   recomputa el verificador \qquad
   \tikz\node[herm,minimum width=4mm,minimum height=2.6mm]{}; hermanos:
   \emph{la prueba} $\pi$};
\end{tikzpicture}
\caption{Árbol de Merkle sobre ocho transacciones (doble SHA-256; se muestran los
primeros 6 dígitos hexadecimales, calculados realmente). Para convencer a
cualquiera de que $T_3$ está adentro basta con enviar tres hashes hermanos, no el
bloque entero.}
\label{fig:merkle}
\end{figure}

Lo elegante es la prueba: para convencerte de que tu transacción está adentro,
sólo hace falta darte los hashes hermanos del camino hasta la raíz. Con
$1\,048\,576$ transacciones son \textbf{20 hashes: 640 bytes} en lugar de
gigabytes. El ahorro es logarítmico, y por eso la verificación simplificada de
pagos (SPV) es posible.

\begin{nota}{Seguridad}
Falsificar una prueba implica hallar una colisión de $\dsha$ en algún nivel. La
seguridad del árbol se reduce a la resistencia a colisión de la función base: es
el mismo argumento que el del teorema de Damgård--Merkle.
\end{nota}

\begin{aviso}{La duplicación de la última hoja}
Cuando un nivel tiene un número impar de nodos, el último se duplica. Ésa es la
fuente de la vulnerabilidad \textbf{CVE-2012-2459} de Bitcoin: dos listas de
transacciones distintas podían producir la misma raíz. Buen ejercicio de clase:
explicar por qué la duplicación rompe la inyectividad
(véase el ejercicio~\ref{ej:cve} del apéndice).
\end{aviso}

\begin{laboratorio}{Laboratorio 7.1 — Constructor de árbol y prueba de inclusión}
Se construye el árbol a partir de una lista de transacciones, se pide la prueba
de inclusión de una de ellas y después se altera esa misma transacción.

\textbf{Qué observar.} La cascada: un solo bit distinto en la hoja cambia todos
los nodos del camino hasta la raíz, y la prueba deja de verificar.
\end{laboratorio}

% ===========================================================================
\section{Compromisos, prueba de trabajo y la cadena}
% ===========================================================================

Un compromiso es un sobre cerrado: pones dentro tu decisión, entregas el sobre, y
más tarde lo abres. Quien lo tiene no puede leerlo (\emph{ocultamiento}) y tú ya
no puedes cambiar lo que hay dentro (\emph{vinculación}).

\begin{definicion}[Esquema de compromiso basado en hash]
\textbf{Compromiso.} El emisor elige $m$ y $r \leftarrow \bits^{\lambda}$
aleatorio, calcula $c = H(m \concat r)$ y publica $c$.

\textbf{Apertura.} Revela $(m, r)$; el verificador recomputa y compara con $c$.

\smallskip
\emph{Ocultamiento}: $c$ no revela información sobre $m$; se apoya en la
resistencia a preimagen y en que $r$ tenga suficiente entropía.

\emph{Vinculación}: hallar $(m', r') \neq (m, r)$ con el mismo $c$ equivale a una
colisión de $H$.
\end{definicion}

\begin{figure}[htbp]
\centering
\begin{tikzpicture}[
  font=\footnotesize,
  act/.style={draw=inkc!55,fill=inkc!8,rounded corners=2pt,
              minimum width=18mm,minimum height=6mm},
  msg/.style={draw=violetdeep,fill=violetc!12,rounded corners=2pt,inner sep=3pt,
              font=\scriptsize},
  >={Stealth[length=2mm]}
]
\node[act] (A) at (0,0) {Alice};
\node[act] (B) at (8.4,0) {Bob};
\draw[inkc!35,dashed] (A) -- ++(0,-6.4);
\draw[inkc!35,dashed] (B) -- ++(0,-6.4);

% fase 1
\node[msg,anchor=west] at (0.35,-1.0)
  {elige $m \in \{\text{águila},\text{sol}\}$ y sortea $r$};
\draw[->,violetdeep,thick] (0,-2.1) -- node[above,font=\scriptsize,violetdeep]
  {$c = H(m \concat r)$} (8.4,-2.1);
\node[font=\scriptsize,coralc,anchor=north west] at (8.7,-2.2)
  {no puede leer $m$};

% fase 2
\draw[<-,amberc,thick] (0,-3.6) -- node[above,font=\scriptsize,amberc]
  {Bob anuncia su apuesta} (8.4,-3.6);

% fase 3
\draw[->,mintc,thick] (0,-5.0) -- node[above,font=\scriptsize,mintc]
  {apertura: $(m, r)$} (8.4,-5.0);
\node[font=\scriptsize,mintc,anchor=north west] at (8.7,-5.1)
  {verifica $H(m\concat r) = c$};

% llaves de fase
\draw[decorate,decoration={brace,amplitude=4pt},inkc!55]
  (-1.35,-1.8) -- (-1.35,-2.4)
  node[midway,left=3pt,font=\scriptsize,inkc]{fase 1};
\draw[decorate,decoration={brace,amplitude=4pt},inkc!55]
  (-1.35,-4.7) -- (-1.35,-5.3)
  node[midway,left=3pt,font=\scriptsize,inkc]{fase 2};

\node[anchor=north west,font=\scriptsize,inkc,align=left,text width=118mm]
  at (-1.35,-6.0)
  {\textbf{Vinculación:} tras la fase 1, Alice ya no puede cambiar $m$ sin
   contradecir $c$. \textbf{Ocultamiento:} hasta la fase 2, $c$ no le dijo nada
   a Bob ---\emph{siempre que $r$ exista}.};
\end{tikzpicture}
\caption{El volado por teléfono. Dos desconocidos deciden algo a distancia sin
confiar el uno en el otro, con un hash como único árbitro.}
\label{fig:compromiso}
\end{figure}

\begin{aviso}{Sin $r$ no hay ocultamiento}
Cuando el espacio de $m$ es pequeño, el adversario simplemente hashea todas las
opciones. Con dos opciones ---águila o sol--- le toma dos hashes. El valor
aleatorio $r$ no es un adorno: es lo que vuelve impracticable esa búsqueda.
\end{aviso}

\subsection{Minar es resolver una preimagen parcial}

La prueba de trabajo pide hallar un \emph{nonce} tal que
\[
  H(\text{bloque} \concat \text{nonce}) < \text{objetivo}.
\]
Es el único problema difícil de esta lista que se resuelve por fuerza bruta
\emph{a propósito}: la dificultad es el producto, no el defecto. Y como cada
bloque incluye el hash del anterior, alterar uno invalida todos los siguientes.

\begin{figure}[htbp]
\centering
\begin{tikzpicture}[
  font=\scriptsize,
  bloque/.style={draw=mintc,fill=mintc!8,rounded corners=3pt,
                 minimum width=36mm,minimum height=27mm,anchor=north west},
  roto/.style={draw=coralc,fill=coralc!8,rounded corners=3pt,
               minimum width=36mm,minimum height=27mm,anchor=north west},
  campo/.style={font=\scriptsize\ttfamily,anchor=north west,inkc},
  >={Stealth[length=2.2mm]}
]
\foreach \i/\est/\prev/\nonce/\hash in {%
  0/bloque/000000000000/10968/0009858c4cf0,
  1/bloque/0009858c4cf0/11581/0003e22f1034,
  2/bloque/0003e22f1034/11945/000cec048d92}{
  \pgfmathtruncatemacro{\num}{\i+1}
  \node[\est] (B\i) at ({\i*4.4},0) {};
  \node[anchor=north west,font=\bfseries\small,mintc]
    at ($(B\i.north west)+(2mm,-1.5mm)$) {Bloque \num};
  \node[campo] at ($(B\i.north west)+(2mm,-7mm)$)  {prev: \prev};
  \node[campo] at ($(B\i.north west)+(2mm,-11mm)$) {nonce: \nonce};
  \node[campo] at ($(B\i.north west)+(2mm,-15mm)$) {hash:};
  \node[anchor=north west,font=\scriptsize\ttfamily,mintc!60!black]
    at ($(B\i.north west)+(11mm,-15mm)$) {\hash};
  \node[anchor=north west,font=\scriptsize,inkc!70]
    at ($(B\i.north west)+(2mm,-20.5mm)$) {$\approx 11\,000$ intentos};
}
\draw[->,mintc,thick] (B0.east|-B0.center) -- (B1.west|-B1.center);
\draw[->,mintc,thick] (B1.east|-B1.center) -- (B2.west|-B2.center);

\node[anchor=north west,align=left,font=\scriptsize,coralc,text width=125mm]
  at (0,-3.2)
  {\textbf{La cascada.} Edita el contenido del bloque~1 y su hash deja de empezar
   con \texttt{000}: el bloque~1 se invalida, el campo \texttt{prev} del
   bloque~2 ya no coincide, y el 3 tampoco. Para reescribir la historia hay que
   volver a minar ese bloque \emph{y todos los posteriores}. Eso es el costo de
   mentir.};
\end{tikzpicture}
\caption{Tres bloques minados de verdad con dificultad de 3 ceros hexadecimales.
Cada cero adicional multiplica por 16 el trabajo esperado; la red de Bitcoin
opera hoy alrededor de 19 ceros.}
\label{fig:cadena}
\end{figure}

\begin{nota}{Por qué SHA-256 se ejecuta en masa}
Un equipo comercial de minado ronda los 580\,TH/s: en un segundo calcula lo que
un núcleo de tu procesador tardaría más de medio año. SHA-256 es el sueño de un
diseñador de circuitos ---64 rondas fijas, sin memoria ni bifurcaciones--- y por
eso un ASIC entrega un hash por ciclo de reloj. Las funciones \emph{memory-hard}
como Argon2 o RandomX invierten el argumento: si el cuello de botella es la
memoria y no la aritmética, el fabricante de silicio pierde su ventaja, porque la
DRAM se la vende el mismo proveedor a todo el mundo.
\end{nota}

\begin{laboratorio}{Laboratorios 8.1 y 8.2 — Volado por teléfono; prueba de trabajo}
En el primero se ejecuta el protocolo de compromiso y luego se repite quitando
$r$: Bob rompe el ocultamiento al instante. En el segundo se minan tres bloques y
se edita el contenido de uno para ver la cascada de invalidación.

\textbf{Qué observar.} Cada cero hexadecimal adicional multiplica por 16 el
trabajo esperado. Por eso el ajuste de dificultad mide trabajo real y no buena
suerte.
\end{laboratorio}

% ===========================================================================
\section{Síntesis: qué se rompe exactamente cuando falla cada propiedad}
% ===========================================================================

Hasta aquí las propiedades se enunciaron por separado. Conviene ahora invertir la
pregunta: si mañana alguien rompe la resistencia a colisión de la función que
estás usando, ¿qué deja de funcionar y qué sigue en pie? La respuesta no es
«todo».

\begin{caso}{El caso canónico: la firma digital}
Nadie firma un documento: \textbf{se firma su hash}. Si el atacante consigue dos
textos con la misma huella ---uno inocuo y otro fraudulento--- basta con hacer
firmar el primero para quedarse con una firma válida sobre el segundo. La clave
privada nunca se compromete y, aun así, la firma es falsificable. Ésa es la
consecuencia exacta de perder la resistencia a colisión, y es lo que ocurrió con
el certificado de CA falso construido sobre MD5 en 2008 y con el malware Flame en
2012.
\end{caso}

\begin{table}[htbp]
\centering\small
\caption{Mapa de consecuencias. Una propiedad rota no derriba el sistema entero:
derriba los elementos cuya seguridad se reduce a ella.}
\label{tab:sintesis}
\begin{tabularx}{\textwidth}{@{}>{\raggedright\arraybackslash}p{24mm}
                               >{\raggedright\arraybackslash}X
                               >{\raggedright\arraybackslash}p{34mm} c@{}}
\toprule
\textbf{Propiedad} & \textbf{Si falla, el adversario puede\ldots} &
\textbf{Dónde duele primero} & \textbf{Costo ideal}\\
\midrule
\textbf{Preimagen} \newline {\scriptsize one-way} &
  Invertir huellas publicadas: recuperar contraseñas de un volcado, leer el
  contenido de un compromiso antes de que lo abran. &
  Ocultamiento del compromiso · contraseñas · direcciones &
  $2^{n}$\\
\addlinespace
\textbf{2.\textsuperscript{a} preimagen} \newline {\scriptsize target collision} &
  Sustituir un objeto ya publicado y aceptado por otro sin alterar ningún hash
  previamente difundido. &
  Encadenamiento de bloques · prueba de Merkle · documento firmado &
  $2^{n}$\\
\addlinespace
\textbf{Colisión} \newline {\scriptsize collision-free} &
  Fabricar de antemano dos objetos intercambiables: hacer firmar uno y exhibir el
  otro, o abrir un mismo compromiso de dos maneras. &
  Firma digital · vinculación del compromiso · raíz de Merkle &
  $2^{n/2}$\\
\addlinespace
\textbf{Colisión de prefijo elegido} \newline {\scriptsize chosen-prefix} &
  Dar a las dos versiones un encabezado realista fijado de antemano ---nombres,
  montos, campos de un certificado X.509---, no dos archivos basura. &
  Documentos con formato · certificados · PGP &
  $\ge 2^{n/2}$\\
\addlinespace
\textbf{Inmunidad a extensión de longitud} \newline
  {\scriptsize propiedad de la construcción} &
  Forjar etiquetas de autenticación para mensajes que nunca vio, sin conocer la
  clave, si el MAC es el ingenuo $H(k \concat m)$. &
  MAC artesanal · APIs firmadas · BIP32 &
  ---\\
\addlinespace
\textbf{Salida indistinguible de aleatoria} \newline
  {\scriptsize avalancha · uniformidad} &
  Sesgar o anticipar valores derivados del hash: el nonce ganador, el resultado
  de un sorteo, la clave hija de un monedero jerárquico. &
  Prueba de trabajo · aleatoriedad pública · HMAC &
  ---\\
\addlinespace
\textbf{Costo de evaluación controlable} \newline {\scriptsize sólo para KDF} &
  Evaluar miles de millones de candidatos por segundo en GPU o ASIC: la
  velocidad, virtud en todo lo demás, aquí es el defecto. &
  Almacenamiento de contraseñas · cifrado de monederos &
  ---\\
\bottomrule
\end{tabularx}
\end{table}

\begin{figure}[htbp]
\centering
\begin{tikzpicture}[
  font=\small,
  prop/.style={draw=#1,fill=#1!10,rounded corners=3pt,minimum width=34mm,
               minimum height=9mm,font=\small\bfseries,text=#1},
  obj/.style={draw=inkc!40,fill=inkc!5,rounded corners=3pt,minimum width=44mm,
              minimum height=9mm,font=\small,align=center},
  >={Stealth[length=2.4mm]}
]
\node[prop=skyc]   (p1) at (0,1.6)  {Preimagen};
\node[prop=amberc] (p2) at (0,0)    {2.\textsuperscript{a} preimagen};
\node[prop=coralc] (p3) at (0,-1.6) {Colisión};

\node[obj] (o1) at (7,1.6)  {lo que debe permanecer\\\textbf{oculto}};
\node[obj] (o2) at (7,0)    {lo que ya fue\\\textbf{publicado y aceptado}};
\node[obj] (o3) at (7,-1.6) {lo que será elegido\\\textbf{por el adversario}};

\draw[->,skyc,thick]   (p1) -- (o1);
\draw[->,amberc,thick] (p2) -- (o2);
\draw[->,coralc,thick] (p3) -- (o3);
\node[anchor=north,font=\footnotesize\itshape,inkc] at (3.5,-2.4)
  {\ldots protege\ldots};
\end{tikzpicture}
\caption{La regla que ordena la tabla~\ref{tab:sintesis}: cada resistencia
custodia una clase distinta de objeto.}
\label{fig:regla}
\end{figure}

\subsection{Una nota sobre el adversario cuántico}

El algoritmo de Grover reduce la búsqueda de preimagen de $2^{n}$ a
$\approx 2^{n/2}$, y el de Brassard--Høyer--Tapp llevaría las colisiones a
$2^{n/3}$ ---pero exigiendo memoria cuántica de acceso aleatorio en una escala
que hoy nadie considera realista. Bajo modelos con memoria acotada, la ventaja
cuántica sobre colisiones casi desaparece.

La consecuencia práctica es sencilla: SHA-256 conserva del orden de 128 bits
frente a preimagen cuántica y \textbf{no es la pieza urgente}. La presión
post-cuántica está en la firma ---ECDSA cae ante Shor---, no en la función hash.
De hecho, los esquemas de firma construidos exclusivamente con hash (FIPS 205,
SLH-DSA) son la respuesta, no el problema.

\begin{proposicion}[hash-then-sign]
Sea $\Sigma$ un esquema de firma seguro frente a falsificación existencial, y
$\Sigma^{H}$ el esquema que firma $H(m)$ en lugar de $m$. Para todo falsificador
$\mathcal{A}$ existen $\mathcal{B}$, $\mathcal{C}$ eficientes tales que
\[
  \Adv^{\text{euf-cma}}_{\Sigma^{H}}(\mathcal{A}) \;\le\;
  \Adv^{\text{euf-cma}}_{\Sigma}(\mathcal{B}) \;+\;
  \Adv^{\text{col}}_{H}(\mathcal{C}).
\]
\end{proposicion}

\begin{observacion}
La lectura importante es la contrapuesta. La desigualdad sólo garantiza seguridad
mientras el segundo término sea despreciable; en cuanto la ventaja de colisión
deja de serlo, la cota se vacía y la falsificación es efectiva. Eso es
exactamente lo que ocurrió con MD5 y con SHA-1.
\end{observacion}

\clearpage
% ===========================================================================
\appendix
\section{Resumen de cotas}
% ===========================================================================

\begin{table}[htbp]
\centering\small
\caption{Todo el módulo en nueve renglones.}
\begin{tabular}{@{}l l l@{}}
\toprule
\textbf{Cantidad} & \textbf{Expresión} & \textbf{Para $n = 256$}\\
\midrule
Preimagen (clásica)        & $2^{n}$                  & $2^{256}$\\
Segunda preimagen          & $2^{n}$                  & $2^{256}$\\
Colisión (cumpleaños)      & $\approx 1.1774\cdot 2^{n/2}$ & $\approx 2^{128}$\\
Preimagen (Grover)         & $\approx 2^{n/2}$        & $\approx 2^{128}$\\
Colisión (BHT, RAM cuántica) & $\approx 2^{n/3}$      & $\approx 2^{85}$\\
Avalancha esperada         & $\mathbb{E}[\dH] = n/2$  & $128$ bits\\
Desviación de la avalancha & $\sigma = \sqrt{n}/2$    & $8$ bits\\
Tamaño de prueba de Merkle & $\lceil\log_2 m\rceil$ hashes & 20 hashes para $m = 2^{20}$\\
Costo del diccionario con sal & $u\cdot|D|\cdot c$    & lineal en usuarios\\
\bottomrule
\end{tabular}
\end{table}

% ===========================================================================
\section{Ejercicios propuestos}
% ===========================================================================

\begin{enumerate}[leftmargin=2.4em,itemsep=0.5em]
  \item Demuestra que la función $h$ de la Proposición~\ref{prop:reciproco} es
    también resistente a segunda preimagen para entradas de longitud
    $\neq n$, y explica por qué eso no contradice nada de lo dicho.

  \item\label{ej:cve} \textbf{CVE-2012-2459.} Sea un árbol de Merkle sobre la
    lista $(T_1, T_2, T_3)$, que duplica la última hoja para completar el nivel.
    Exhibe una lista distinta de cuatro transacciones con exactamente la misma
    raíz, y argumenta por qué el atacante no necesita ninguna colisión de
    SHA-256 para construirla.

  \item Un servidor autentica con $t = H(m \concat k)$ ---la clave \emph{al
    final}---. ¿Sigue funcionando el ataque de extensión de longitud? ¿Qué otro
    problema aparece en su lugar cuando $H$ pierde la resistencia a colisión?

  \item Usando la ecuación~\eqref{eq:cumple}, calcula cuántas direcciones
    distintas hay que generar para que la probabilidad de que dos coincidan
    alcance el 1\,\% con $n = 160$ bits. Compara con el número de direcciones que
    ha visto la red de Bitcoin en toda su historia.

  \item Un sistema guarda $v = H(s \concat p)$ con sal de 128 bits pero usa
    SHA-256 de una sola pasada. Cuantifica, con la figura~\ref{fig:contrasenas},
    cuánto gana el defensor al migrar a Argon2id con $c = 10^{5}$, y explica por
    qué la sal por sí sola no bastaba.

  \item Demuestra que un compromiso $c = H(m)$ \emph{sin} valor aleatorio sigue
    siendo vinculante pero pierde por completo el ocultamiento cuando
    $|\mathcal{M}| = 2$. ¿Cuántos hashes necesita el adversario?
\end{enumerate}

\vfill
\begin{center}
\color{inkc}\small
\rule{0.35\textwidth}{0.4pt}\\[0.6em]
\textsc{ChainNotes} · Módulo 01 de 9 · Siguiente: esquemas de compromiso\\
{\footnotesize Versión LaTeX del apunte interactivo \texttt{temas/hash.html}.
Todas las figuras se generan con TikZ y pgfplots.}
\end{center}

\end{document}
